% sec:rk-butcher — Explicit Runge-Kutta Methods and Butcher Tableaux
\section{Explicit Runge-Kutta Methods and Butcher Tableaux}
\label{sec:rk-butcher}

The finite difference stencils of \cref{ch:finite-differences} convert
spatial derivatives into algebraic operations on grid data.  But a
derivative approximation is not yet a time step: the evolution equations
of numerical relativity and the geodesic equations of the ray tracer both
have the form
%
\begin{equation}
  \frac{\d y}{\d t} = f(t, y) \,, \qquad y(t_0) = y_0 \,,
  \label{eq:ivp}
\end{equation}
%
where $y$ may be a scalar, a vector, or the full state of twenty-four
BSSN fields across a three-dimensional grid.  The task is to advance
$y$ from time $t_n$ to $t_{n+1} = t_n + h$ using only evaluations
of the right-hand side $f$.
\physnote{For geodesic integration, $t$ is the affine parameter
$\lambda$ and $y = (x^\mu, p_\mu)$ is the eight-component
phase-space state.  For the BSSN system, $t$ is coordinate time and $y$
comprises all evolved fields at every grid point.}

\paragraph{From Euler to multi-stage methods.}

The simplest approach evaluates $f$ once at the current state and
extrapolates linearly:
%
\begin{equation}
  y_{n+1} = y_n + h\,f(t_n,\, y_n) \,.
  \label{eq:euler}
\end{equation}
%
This is Euler's method.  A Taylor expansion shows that the local
truncation error --- the error introduced in a single step, assuming
exact input --- is $\order{h^2}$, giving a global error that is first
order in $h$.  For the geodesic integrator, first-order accuracy is
disastrous: tracing a photon orbit around a black hole for thousands of
steps accumulates error so rapidly that the trajectory bears no
resemblance to the true geodesic after a few orbital periods.
\histnote{Euler published this method in 1768; Runge and Kutta
independently developed multi-stage extensions in 1895 and 1901.}

A right-hand side that curves within the interval cannot be captured by a
single slope sample at its start --- but evaluating the slope partway
through the step and correcting the estimate can cancel leading error
terms, much as the midpoint rule for quadrature outperforms the
left-endpoint rule.  The remedy is to sample $f$ at several points within
the interval $[t_n,\, t_{n+1}]$ and combine the samples into a more
accurate estimate of the step.  Each sample is called a \emph{stage}.
With $s$ stages, the method evaluates $f$ exactly $s$ times per step,
and the coefficients that specify where to sample and how to weight
the results constitute the \emph{Runge--Kutta method}.

\paragraph{The general $s$-stage explicit method.}

An explicit Runge--Kutta method with $s$ stages computes intermediate
slopes $k_1, \ldots, k_s$ and assembles the step as follows.  The
first stage evaluates $f$ at the current point:
%
\begin{equation}
  k_1 = f(t_n,\, y_n) \,.
  \label{eq:rk-k1}
\end{equation}
%
This is the slope at the start of the interval --- exactly what Euler's
method uses as its sole estimate.  Each subsequent stage evaluates $f$
at a point constructed from the previous stages:
%
\begin{equation}
  k_i = f\!\biggl(t_n + c_i\,h,\;\;
    y_n + h \sum_{j=1}^{i-1} a_{ij}\,k_j\biggr) \,,
  \qquad i = 2, \ldots, s \,.
  \label{eq:rk-stages}
\end{equation}
%
The key word is \emph{explicit}: stage $k_i$ depends only on the
stages $k_1, \ldots, k_{i-1}$ already computed, never on $k_i$
itself or any later stage.  This makes the evaluation a straightforward
forward sweep --- no systems of equations need to be solved.

The final update combines all stages with a set of weights:
%
\begin{equation}
  y_{n+1} = y_n + h \sum_{i=1}^{s} b_i\,k_i \,.
  \label{eq:rk-update}
\end{equation}
%
The method is fully specified by three sets of coefficients: the
\emph{nodes} $c_i$ (where in the interval to evaluate), the
\emph{coupling matrix} $A = (a_{ij})$ (how to combine previous stages),
and the \emph{weights} $b_i$ (how to combine all stages into the final
step).

\paragraph{The Butcher tableau.}

John Butcher introduced a compact notation that displays these
coefficients in a single table \cite{HairerNorsettWanner1993}:
%
\begin{equation}
  \renewcommand{\arraystretch}{1.3}
  \begin{array}{c|c}
    \mathbf{c} & A \\[2pt]
    \hline
    & \mathbf{b}^{\!\top}
  \end{array}
  \;\;=\;\;
  \begin{array}{c|cccc}
    c_1    & 0      &        &        & \\
    c_2    & a_{21} & 0      &        & \\
    \vdots & \vdots & \ddots & \ddots & \\
    c_s    & a_{s1} & \cdots & a_{s,s-1} & 0 \\[2pt]
    \hline
           & b_1    & b_2    & \cdots & b_s
  \end{array}
  \label{eq:butcher-tableau}
\end{equation}
%
The upper-triangular zeros in $A$ encode the explicit structure: no
stage depends on itself or any later stage.  Reading the tableau
row by row reconstructs the method: row $i$ of $A$ gives the
coefficients for building the argument of $k_i$, and the bottom
row gives the weights for the final combination.
\warnnote{The nodes should satisfy the \emph{row-sum condition}
$c_i = \sum_{j=1}^{i-1} a_{ij}$ (with $c_1 = 0$).  This condition
is necessary for the method to be consistent, i.e., to reproduce the
exact solution for constant right-hand sides.  Without it, some or
all of the order conditions beyond first order fail, degrading the
method's accuracy.}

\paragraph{The classical fourth-order method.}

The most widely known Runge--Kutta method uses four stages to achieve
fourth-order accuracy.  Its Butcher tableau is
%
\begin{equation}
  \renewcommand{\arraystretch}{1.3}
  \begin{array}{c|cccc}
    0   &     &     &     & \\
    1/2 & 1/2 &     &     & \\
    1/2 & 0   & 1/2 &     & \\
    1   & 0   & 0   & 1   & \\[2pt]
    \hline
        & 1/6 & 1/3 & 1/3 & 1/6
  \end{array}
  \label{eq:rk4-tableau}
\end{equation}
%
The structure is transparent: sample the slope at $t_n$
(beginning), twice at $t_n + h/2$ (midpoint, using different
estimates to get there), and once at $t_n + h$ (end), then
take a weighted average with the Simpson's-rule-like
weights $\frac{1}{6}$, $\frac{1}{3}$, $\frac{1}{3}$,
$\frac{1}{6}$.  Writing out the stages explicitly gives the
familiar formulae:
%
\begin{align}
  k_1 &= f(t_n,\; y_n) \,, \notag \\
  k_2 &= f(t_n + \tfrac{h}{2},\; y_n + \tfrac{h}{2}\,k_1) \,, \notag \\
  k_3 &= f(t_n + \tfrac{h}{2},\; y_n + \tfrac{h}{2}\,k_2) \,,
  \label{eq:rk4-stages} \\
  k_4 &= f(t_n + h,\; y_n + h\,k_3) \,, \notag \\[6pt]
  y_{n+1} &= y_n + \frac{h}{6}\bigl(k_1 + 2\,k_2 + 2\,k_3 + k_4\bigr) \,.
  \notag
\end{align}
%
Each stage is a single right-hand-side evaluation, and the final
combination requires only addition and scalar multiplication.
For the geodesic integrator, each evaluation of $f$ computes the
Hamiltonian right-hand side $(\dot{x}^\mu, \dot{p}_\mu)$ from the
current phase-space state --- the dominant cost per step.

\paragraph{Order and the cost of accuracy.}

A method has \emph{order} $p$ if the local truncation error satisfies
$\norm{y(t_{n+1}) - y_{n+1}} = \order{h^{p+1}}$ for exact input
$y_n = y(t_n)$, where $y(t)$ is the exact solution of
\cref{eq:ivp}.  The global error accumulated over a fixed time
interval is then $\order{h^p}$, provided $f$ is Lipschitz continuous
in $y$: doubling the resolution halves the step size and reduces the
error by a factor of $2^p$.
For the classical RK4, $p = 4$, so doubling the resolution reduces
the global error by a factor of $16$ --- the same convergence
diagnostic used for the spatial stencils in
\cref{sec:fd-convergence}.
\xrefnote{The self-convergence test of \cref{sec:fd-convergence}
applies unchanged to ODE integrators: run at three resolutions with
step-size ratio $4:2:1$ and verify $\rho = 2^p$.}

Achieving order $p$ with $s$ stages requires that the Taylor
expansion of $y_{n+1}$ match the Taylor expansion of the exact
solution through $\order{h^p}$.  These \emph{order conditions}
constrain the Butcher tableau coefficients.  Matching the Taylor
expansion to order $p$ generates a set of nonlinear algebraic
conditions on the entries of $A$ and $\mathbf{b}$, and the number
of these conditions grows combinatorially with $p$ --- at fifth order
there are seventeen conditions but only the fifteen free parameters
of a five-stage method.  The sharp lower bounds on the stage count,
established through work by Butcher and others
\cite{HairerNorsettWanner1993}, are:
%
\begin{equation}
  \renewcommand{\arraystretch}{1.2}
  \begin{array}{l|cccccc}
    \text{order}\; p & 1 & 2 & 3 & 4 & 5 & 6 \\
    \hline
    \text{min.\ stages} & 1 & 2 & 3 & 4 & 6 & 7
  \end{array}
  \label{eq:butcher-barrier}
\end{equation}
%
For orders one through four, the minimum stage count equals $p$: the
classical RK4 is optimal in this sense.  But fifth-order accuracy
demands at least six stages --- one more than the naive expectation.
This \emph{Butcher barrier} is the reason that the Dormand--Prince
method of the next section uses seven stages rather than five: the
extra stages buy fifth-order accuracy and, crucially, a built-in error
estimate that enables adaptive step-size control.
\implnote{The codebase's \texttt{dormandPrinceStep} function in
\codemodule{Geodesic.Solver} evaluates seven stages per step,
directly mirroring the Butcher tableau structure of
\cref{eq:rk-stages,eq:rk-update}.}

\paragraph{From tableau to code.}

The tableau cleanly separates the \emph{method} --- encoded in
the coefficients $(A, \mathbf{b}, \mathbf{c})$ --- from the
\emph{algorithm} that evaluates stages and accumulates the result.  This separation is exactly what appears in the codebase:
changing from one Runge--Kutta method to another means changing
the tableau coefficients, not rewriting the integration loop.
\Cref{sec:rk-dormand-prince} exploits this structure to introduce
the Dormand--Prince embedded pair, where two sets of weights ---
$\mathbf{b}$ for a fifth-order solution and $\hat{\mathbf{b}}$ for a
fourth-order companion --- share the same stages, providing a local
error estimate at no extra cost.
